

\documentclass[11pt,a4paper]{article} 

%------------------------------------------------------------------------------
%	REQUIRED PACKAGES AND  CONFIGURATIONS
%------------------------------------------------------------------------------
% PACKAGES FOR TITLES
\usepackage{titlesec}
\usepackage{color}

% PACKAGES FOR LANGUAGE AND FONT
\usepackage[utf8]{inputenc}
\usepackage[spanish]{babel}
\usepackage[T1]{fontenc} % Font encoding

% PACKAGES FOR IMAGES
\usepackage{graphicx}
\graphicspath{{Images/}}
\usepackage{eso-pic} % For the background picture on the title page
\usepackage{subfig} % Numbered and caption subfigures using \subfloat
\usepackage{caption} % Coloured captions
\usepackage{transparent}

% STANDARD MATH PACKAGES
\usepackage{amsmath}
\usepackage{amsthm}
\usepackage{bm}
\usepackage[overload]{empheq}  % For braced-style systems of equations

% PACKAGES FOR TABLES
\usepackage{tabularx}
\usepackage{longtable} % tables that can span several pages
\usepackage{colortbl}

% PACKAGES FOR ALGORITHMS (PSEUDO-CODE)
\usepackage{algorithm}
\usepackage{algorithmic}

% PACKAGES FOR REFERENCES & BIBLIOGRAPHY
\usepackage[colorlinks=true,linkcolor=black,anchorcolor=black,citecolor=black,filecolor=black,menucolor=black,runcolor=black,urlcolor=black]{hyperref} % Adds clickable links at references
\usepackage{cleveref}
\usepackage[square, numbers, sort&compress]{natbib} % Square brackets, citing references with numbers, citations sorted by appearance in the text and compressed
\bibliographystyle{plain} % You may use a different style adapted to your field

% PACKAGES FOR THE APPENDIX
\usepackage{appendix}

% PACKAGES FOR ITEMIZE & ENUMERATES 
\usepackage{enumitem}

% OTHER PACKAGES
\usepackage{amsthm,thmtools,xcolor} % Coloured "Theorem"
\usepackage{comment} % Comment part of code
\usepackage{fancyhdr} % Fancy headers and footers
\usepackage{lipsum} % Insert dummy text
\usepackage{tcolorbox} % Create coloured boxes (e.g. the one for the key-words)
\usepackage{comment} % Create comments

%-------------------------------------------------------------------------
%	NEW COMMANDS DEFINED
%-------------------------------------------------------------------------
% EXAMPLES OF NEW COMMANDS -> here you see how to define new commands
\newcommand{\bea}{\begin{eqnarray}} % Shortcut for equation arrays
\newcommand{\eea}{\end{eqnarray}}
\newcommand{\e}[1]{\times 10^{#1}}  % Powers of 10 notation
\newcommand{\mathbbm}[1]{\text{\usefont{U}{bbm}{m}{n}#1}} % From mathbbm.sty
\newcommand{\pdev}[2]{\frac{\partial#1}{\partial#2}}
% NB: you can also override some existing commands with the keyword \renewcommand

%----------------------------------------------------------------------------
%	ADD YOUR PACKAGES (be careful of package interaction)
%----------------------------------------------------------------------------


%----------------------------------------------------------------------------
%	ADD YOUR DEFINITIONS AND COMMANDS (be careful of existing commands)
%----------------------------------------------------------------------------


% Do not change Configuration_files/config.tex file unless you really know what you are doing. 
% This file ends the configuration procedures (e.g. customizing commands, definition of new commands)
\input{Configuration_files/config}

% Insert here the info that will be displayed into your Title page 
% -> title of your work
\renewcommand{\title}{Estudio de los efectos del sistema immune y la quimioterapia en el crecimiento de tumores avasculares modelado por autómatas celulares}
% -> author name and surnam separated by \\
\newcommand{\authors}{Gómez Rueda, José M.\\ Galindo Álvarez, Julián M.\\ Ramos González, Víctor.\\ Rodriguez Gallego, José A.}

% -> Authors emails. Keep the previous order 
\newcommand{\contacts}{\href{mailto:josemariagomezrueda@gmail.com}{josemariagomezrueda@gmail.com}\\\href{mailto:julian.g.a1996@gmail.com}{julian.g.a1996@gmail.com}\\\href{mailto:ramgonvictor@gmail.com}{ramgonvictor@gmail.com}\\josea@gmail.com}
% -> advisor name and surname
\newcommand{\advisor}{Sancho Caparrini, Fernando }
% IF AND ONLY IF you need to modify the co-supervisors you also have to modify the file Configuration_files/title_page.tex (ONLY where it is marked)
\newcommand{\firstcoadvisor}{Sancho Caparrini, Fernando} % insert if any otherwise comment
\newcommand{\secondcoadvisor}{Hurtado Mendoza, Ignacio} % insert if any otherwise comment
% -> academic year
\newcommand{\YEAR}{2021-2022}

% -> Resources (complement resources: documents, models, etc). Separated by \\
\newcommand{\otherResources}{\href{https://github.com/julgalalv/chemotherapy-effects-SCAM}{@github:MSCAM}}

% -> abstract (only in English)
\renewcommand{\abstract}{En este trabajo se tratan de estudiar los efectos de la terapia en el desarrollo del cáncer avascular sobre la base de un modelo estocástico de autómatas celulares, presentando un modelo basado en la caracterización de tres procesos principales: crecimiento tumoral, acción inmunitaria  y acción farmacológica basada en quimioterapia. Más concretamente, el modelo presentado, MSCAM (\textit{Modified SCAM}), está basado en el modelo SCAM \cite{Pourhasanzade2019} incluyendo, a fin de que el modelo sea más compatible con el comportamiento biológico, un desglose de la actividad inmunitaria (innata y adquirida) junto con los procesos relacionados con ella (reconocimiento, respuesta inmune innata, activación de agentes inmunes específicos, respuesta inmune adquirida), así como la modelización de la respuesta farmacalógica contemplando tanto el factor de farmaco-resistencia (innata e inducida por los fármacos) así como el factor de repoblación tumoral. El trabajo abarca el estudio de distintos aspectos incluyendo análisis de la evolución tumoral e inmunitaria en ausencia quimioterapéutica, bajo la interrupción del tratamiento tras un solo ciclo de quimioterapia y bajo los modos de terapia de administración periódica de fármacos. La fracción de crecimiento, la fracción necrótica y el volumen tumoral se utilizan como marcadores de salida junto a una presentación gráfica de crecimiento en 2 dimensiones. Como validación de los resultados obtenidos se realiza una comparativa con los propuestos en el artículo de referencia \cite{Pourhasanzade2019}, basado en el contraste con datos bibliográficos de estudios in vivo.}

% -> key-words (only in English)
\newcommand{\keywords}{Sthocastic Cellular Automata, Cancer Modeling \& Analysis, Cancer therapies}

%-------------------------------------------------------------------------
%	BEGIN OF YOUR DOCUMENT
%-------------------------------------------------------------------------
\begin{document}

%-----------------------------------------------------------------------------
% TITLE PAGE
%-----------------------------------------------------------------------------
% Do not change Configuration_files/TitlePage.tex (Modify it IF AND ONLY IF you need to add or delete the Co-advisors)
% This file creates the Title Page of the document
% DO NOT REMOVE SPACES BETWEEN LINES!

\AddToShipoutPicture*{\BackgroundPic}

\includegraphics[width=0.6\textwidth]{logo_polimi_ing_indinf.eps}

\vspace{1.3cm}

{\small \bfseries \textsc{ \color{black!60} MODELOS DE COMPUTACIÓN EN VIDA ARTIFICIAL }}\\

\vspace{0.3cm}

{\LARGE \textbf{\color{bluePoli} \title}}\\

\normalfont

\centerline{\rule{1.0\textwidth}{0.4pt}}

\begin{center}
\begin{minipage}[t]{.28\textwidth}

\scriptsize{\textbf{Autores:}} \\
\authors

\vspace{0.2cm}


\scriptsize{\textbf{Contactos:}} \\
\contacts

\vspace{0.2cm}

\scriptsize{\textbf{Profesor:}} \\
\advisor \\
\\
\textbf{Profesores Asig:} \\ % leave it if any co-advisor otherwise comment
\firstcoadvisor \\ % leave it if any co-advisor otherwise comment
\secondcoadvisor \\ % leave it if you have more that one co-advisor otherwise comment (if you have more than two co-advisors just copy&paste this line writing \thirdcoadvisor, \fourthcoadvisor, ecc. (REMEMBER to modify also the main.txt)
\\ % leave it if any co-advisor otherwise comment
\textbf{Curso académico:} \\
\YEAR 

\vspace{0.2cm}


\scriptsize{\textbf{Recursos adicionales:}} \\
\otherResources\\


\end{minipage}
\begin{minipage}[t]{.70\textwidth}
\noindent \textbf{\color{bluePoli} Abstract:} {\abstract}
\end{minipage}
\end{center}

\vspace{15pt}

\begin{tcolorbox}[arc=0pt, boxrule=0pt, colback=bluePoli!60, width=\textwidth, colupper=white]
    \textbf{Keywords:} \keywords
\end{tcolorbox}

\vspace{12pt}

%%%%%%%%%%%%%%%%%%%%%%%%%%%%%%
%%     THESIS MAIN TEXT     %%
%%%%%%%%%%%%%%%%%%%%%%%%%%%%%%

%-----------------------------------------------------------------------------
% INTRODUCTION
%-----------------------------------------------------------------------------
\section{Introducción}
\label{sec:introduction}
\input{material/introduction.tex}

%-----------------------------------------------------------------------------
% Materials and Methods : MSCAM MODEL
%-----------------------------------------------------------------------------
\section{Modelo MSCAM}
\label{sec:model}
\input{material/model.tex}

%-----------------------------------------------------------------------------
% IMPLEMENTATION
%-----------------------------------------------------------------------------
%Optional
\section{Implementación}
\label{sec:impl}
\input{material/impl.tex}

%-----------------------------------------------------------------------------
% RESULTS AND DISCUSSION
%-----------------------------------------------------------------------------
%Optional
\section{Resultados y discusión}
\label{sec:results}
\input{material/results.tex}

%-----------------------------------------------------------------------------
% CONCLUSION
%-----------------------------------------------------------------------------
\section{Resultados y discusión}
\label{sec:conclusions}
\input{material/conclusions.tex}

\section{Bibliography and citations}



%-----------------------------------------------------------------------------
% BIBLIOGRAPHY
%-----------------------------------------------------------------------------
\bibliographystyle{ieeetran}
\bibliography{bibliography.bib}

%-----------------------------------------------------------------------------
% APPENDICES
%-----------------------------------------------------------------------------

%\newpage
%
%\appendix
%\input{material/appendices.tex}

%-------------------------------------------------------------------------
%	END OF YOUR DOCUMENT
%-------------------------------------------------------------------------
\end{document}